%% One-page Honours Project Summary Template (N Hale 09/12/2025)

\documentclass[12pt,a4paper]{article}

\usepackage[utf8]{inputenc}
\usepackage[T1]{fontenc}

% Clean, reliable margins
\usepackage[a4paper,margin=25mm]{geometry}

\usepackage{xcolor}
\usepackage{hyperref}
\hypersetup{
	colorlinks=true,
	linkcolor=blue,
	filecolor=blue,
	citecolor=blue,
	urlcolor=blue
}

\usepackage{url}
\urlstyle{same}

\usepackage{amsfonts}
\usepackage{amsmath}
\usepackage{graphicx}
\usepackage[expansion=false]{microtype}
\usepackage{lipsum}

% If you aren't using natbib commands, you can remove this.
\usepackage[sort,numbers]{natbib}
\renewcommand{\bibsection}{}

\usepackage{enumitem}
\setlist[itemize]{nosep,leftmargin=*}

% For non-floating captions (keeps figure on the page)
\usepackage{caption}

% Paragraph behaviour
\setlength{\parindent}{0pt}
\setlength{\parskip}{2mm plus 1pt}



\usepackage{titling}

% Pull the whole title block upward
\setlength{\droptitle}{-3.5cm}

% Tighten the internal spacing too (optional but nice)
\pretitle{\begin{center}\Large\bfseries}
	\posttitle{\par\end{center}\vspace{-1.0em}}
\preauthor{\begin{center}\normalsize}
	\postauthor{\par\end{center}\vspace{-2.0em}}
\predate{\begin{center}\normalsize}
	\postdate{\par\end{center}\vspace{-4.0em}}

% -------------------------------------------------------
% Title info
% -------------------------------------------------------
\title{\textbf{Planning as Autoregressive Sequence Modelling: Decoding Strategies}\\[-.75em]}
\author{Dian\ Kruger\\[-0.8em]
	\hspace*{0cm}\\
	Supervisor: Dr S.\ Josias}
\date{}

\begin{document}
	
	\maketitle
	% No page numbers
	\thispagestyle{empty}
	
	% -------------------------------------------------------
	% One-page structure
	% -------------------------------------------------------
	
	The problem of finding a sequence of actions that achieves a goal in a given environment is known as 
	planning. This project investigates the planning problem of finding sequences of moves, called a trajectories,
	from a starting position to a goal position in different grid-based
	mazes. Figure~\ref{fig:fig1} shows an example of such a trajectory.
	A transformer is a type of neural network
	commonly used in language models to predict probabilities of sequences of words.
	This project, however, trains a Transformer to predict probabilities of sequences of actions.
	The rule that is used to choose the next move in a trajectory from these probabilities
	is known as the decoding strategy. 
	The central research question for this project then is: how does the choice of decoding strategy affect 
	the quality of the generated trajectories for maze planning?

	\begin{figure}[htbp]
		\centering
		\captionsetup{font=footnotesize,skip=3pt}
		\begin{minipage}[t]{0.36\linewidth}
			\centering
			\includegraphics[width=0.61\linewidth]{../example_trajectory.png}
			\caption{Example of a valid trajectory.}
			\label{fig:fig1}
		\end{minipage}\hfill
		\begin{minipage}[t]{0.61\linewidth}
			\centering
			\includegraphics[width=0.94\linewidth]{../partial_observation_vs_true_state_V2.png}
			\caption{The complete maze state (left) and the information available to the model in the partially observed setting (right).}
			\label{fig:partial_observation}
		\end{minipage}
	\end{figure}

	A series of experiments was conducted to compare three decoding strategies: greedy 
	decoding, beam search with width $B=2$, and beam search with width $B=3$. They were compared
	in terms of how often they reached the goal, how many steps they made until completion and
	how much computation time they required. 
	Greedy decoding always chooses the move with the highest predicted probability.
	Beam search keeps track of
	the $B$ most probable sequences of moves at each step, extending by one move
	and choosing the most probable out of the at $B$ retained sequences at the end~\cite{jurafsky2026speech}.

	The first experiment was conducted in maze environments where the model was given 
	the entire maze layout and the start and goal positions.
	The model therefore knows its exact position in the maze at every step.
	Beam search reached the goal more often and took fewer moves than greedy decoding,
	with the wider beam performing best. Although each move took slightly longer to compute,
	beam search required less computation time per planning attempt overall.

	The second experiment compared the same decoding strategies, but used a larger training dataset
	consisting of trajectories which 
	the model was not given the complete maze layout or the coordinates of its position and the goal. 
	Instead, at each position, it received a $3\times3$ view of the surrounding squares, the goal’s 
	position relative to its current position, and the previous observations and moves. 
	At every position in the maze, the model only "sees" the information shown in 
	Figure~\ref{fig:partial_observation}.  
	In this setting, beam search reached the goal more often and in fewer moves than
	greedy decoding, with the wider beam performing slightly better. Greedy decoding,
	however, was much faster both per move and per planning attempt.

	\begin{thebibliography}{2}
		\bibitem{jurafsky2026speech}
		Daniel Jurafsky and James H. Martin.\ \textit{Speech and Language Processing}, 3rd edition, online manuscript, 2026.
	\end{thebibliography}


	
	
	
	
	
	% These probability for each move
	% be thought of as: to what degree does the sequence of moves resemble those found in the trainng dataset 
	% if the sequence of moves is extended by this next move. When such a model is trained, it can generate
	% entire sequences of moves, by determining the probabilities of possible first move, chosing one of them
	% according to some rule, plug that move as input into the same model, repeat the previous process to 
	% decide on a second move, plug the first two moves into the mode to determine the third, and repeat
	% untill the goal is reached. The process of generating a sequence of moves, by repeatedly extending the 
	% sequence by one move at a time, is kown as autoregressive sequence modelling, and the rule used to
	% decide what move to take given the predicted probabilities is known as the decoding strategy. This
	% is the focus of the research project: how does the choice of decoding strategy affect the quality
	% of the generated trajectories under autoregressive sequence modelling of maze planning?
	
	

	

\end{document}
